\documentclass[11pt]{article}
\usepackage[margin=1in]{geometry}
\usepackage{amsmath,amssymb,amsthm}
\usepackage[hidelinks]{hyperref}
\newtheorem{theorem}{Theorem}
\newtheorem{lemma}{Lemma}
\newcommand{\fall}[2]{(#1)_{\underline{#2}}}
\title{The odd-degree derivative-root mean identity\\
\large A solution to MathDB 336466}
\author{Principia Math In Progress}
\date{18 September 2026}
\begin{document}
\maketitle
\begin{abstract}
We prove the conjectured alternating binomial relation between the
means of an odd-degree polynomial evaluated at the roots of its
successive derivatives. The proof interpolates the means as a
polynomial in the degree of the derivative. Its otherwise missing
endpoint is an alternating coefficient pairing, which vanishes in
odd degree. The argument includes arbitrary complex coefficients
and roots counted with multiplicity, and gives an explicit
correction in every even degree.
\end{abstract}

\section*{Task 1: the original problem and its full solution}
Let $f\in\mathbb C[x]$ have degree $D\geq1$. For $0\leq\rho<D$,
define
\[
 M_\rho(f)=\frac1{D-\rho}\sum_{f^{(\rho)}(r)=0}f(r),
\]
where roots are counted with their algebraic multiplicities.
In the notation of Wojnar, Wojnar and Brin \cite{W}, this is
$\varphi_{D,0,\rho}$. Their Conjecture 2 (page 9) asserts that
\begin{equation}\label{odd}
 \sum_{\rho=1}^{D-1}(-1)^\rho\binom D\rho M_\rho(f)=0
 \qquad(D\text{ odd}).
\end{equation}
We prove the following stronger formula.

\begin{theorem}\label{main}
Write $f(x)=a g(x)$, where $a\ne0$ and
$g(x)=\sum_{j=0}^D c_jx^{D-j}$ is monic, so $c_0=1$. Then
\begin{equation}\label{general}
 \sum_{\rho=1}^{D-1}(-1)^\rho\binom D\rho M_\rho(f)
 =a(-1)^{D+1}J_D(c),\qquad
 J_D(c)=\sum_{m=0}^D\frac{(-1)^m c_m c_{D-m}}{\binom Dm}.
\end{equation}
In particular, \eqref{odd} holds for every odd $D$.
\end{theorem}

\subsection*{Polynomial interpolation of root power sums}
For an indeterminate $u$ and $0\leq j\leq D$, set
\[
 A_j(u)=c_j\frac{\fall uj}{\fall Dj},\qquad
 \fall uj=u(u-1)\cdots(u-j+1),\qquad \fall u0=1.
\]
For each integer $1\leq n\leq D$, direct differentiation gives
\begin{equation}\label{derivative}
 g_n(x):=\frac{n!}{D!}g^{(D-n)}(x)
       =\sum_{j=0}^n A_j(n)x^{n-j}.
\end{equation}
This is monic of degree $n$, and $A_j(n)=0$ for $j>n$.
Introduce the formal power series
\[
 B(u,z)=\sum_{j=0}^D A_j(u)z^j,\qquad
 P_m(u)=-m[z^m]\log B(u,z)\quad(1\leq m\leq D).
\]
Here $A_0=1$, so the formal logarithm is well-defined. All
coefficient extractions involve only finite sums; analytic
convergence is irrelevant.

\begin{lemma}\label{powers}
Each $P_m$ is a polynomial in $u$ of degree at most $m$, divisible
by $u$. If $r_1,\ldots,r_n$ are the roots of $g_n$, with
multiplicities, then $P_m(n)=\sum_{i=1}^n r_i^m$.
\end{lemma}
\begin{proof}
For integer $n\geq1$, \eqref{derivative} implies
$B(n,z)=\prod_{i=1}^n(1-r_i z)$. Taking formal logarithms and
using $\log(1-rz)=-\sum_{m\geq1}r^mz^m/m$ proves the power-sum
identity, including repeated roots.

In the coefficient of $z^m$ in
$\log(1+\sum_{j\geq1}A_j(u)z^j)$, each product has indices
whose sum is $m$. Since $\deg_u A_j\leq j$, its degree is at
most $m$. Finally, $A_j(0)=0$ for $j\geq1$, so $B(0,z)=1$
and $P_m(0)=0$.
\end{proof}

Consequently,
\begin{equation}\label{H}
 H(u)=c_D+\sum_{m=1}^D c_{D-m}\frac{P_m(u)}u
\end{equation}
is a polynomial of degree at most $D-1$. The quotients in this
definition are polynomial quotients, so evaluation at $u=0$ is
legitimate. For $1\leq n\leq D$, the lemma gives
\begin{equation}\label{means}
 H(n)=\frac1n\sum_{i=1}^n g(r_i)=M_{D-n}(g).
\end{equation}
In particular, $H(D)=M_0(g)=0$.

\subsection*{The missing endpoint and the finite difference}
At $u=0$, differentiating the formal logarithm yields
\[
 P_m'(0)=-m[z^m]\left.\frac{\partial B/\partial u}{B}\right|_{u=0}
        =-m A_m'(0).
\]
The derivative of $\fall um$ at zero is
$(-1)^{m-1}(m-1)!$. It follows that
\[
 P_m'(0)=\frac{(-1)^m m!c_m}{\fall Dm}
        =\frac{(-1)^m c_m}{\binom Dm}.
\]
Substitution into \eqref{H} proves
\begin{equation}\label{endpoint}
 H(0)=c_D+\sum_{m=1}^D c_{D-m}P_m'(0)=J_D(c).
\end{equation}

\begin{proof}[Proof of Theorem \ref{main}]
The $D$th finite difference of a polynomial of degree at most
$D-1$ is zero. Apply this to $H$:
\[
 0=\sum_{\rho=0}^D(-1)^\rho\binom D\rho H(D-\rho)
  =\sum_{\rho=1}^{D-1}(-1)^\rho\binom D\rho M_\rho(g)
    +(-1)^D J_D(c).
\]
We used $H(D)=0$, \eqref{means}, and \eqref{endpoint}.
Since multiplying by $a\ne0$ leaves derivative roots unchanged
and multiplies every mean by $a$, this proves \eqref{general}.

If $D$ is odd, the terms with indices $m$ and $D-m$ in $J_D$
have equal denominators and opposite signs. There is no unpaired
middle index, so $J_D=0$. This proves \eqref{odd}. For $D=1$
the asserted sum is empty and the same proof applies.
\end{proof}

For example, the cubic relation is $-3M_1+3M_2=0$, and the
quintic relation is $-5M_1+10M_2-10M_3+5M_4=0$.
In degree two, \eqref{general} reads
$-2M_1(f)=a(c_1^2/2-2c_2)$, confirming the sign of the even-degree
correction. No assumption of distinct or real roots was used.

\section*{Task 2: the dimension of all universal mean relations}
For each $D$, classify the vectors $(b_1,\ldots,b_{D-1})$ for
which $\sum b_\rho M_\rho(f)=0$ for every degree-$D$ polynomial.
The source's computations suggest dimension one for even $D>2$
and dimension two for odd $D>3$. Prove or correct these dimension
claims, with all small-degree exceptions stated explicitly.

\section*{Task 3: explicit bases for the relation spaces}
Find closed formulas for a basis of each space in Task 2. In odd
degree, complete the binomial vector proved here to a basis.
Determine whether the coefficients of $H(u)$ provide a uniform
construction, and recover the source's tabulated low-degree
relations as checks.

\section*{Task 4: other polynomials evaluated at derivative roots}
Fix $D$ and $0\leq s\leq D$. Replace the value $f(r)$ by an
arbitrary polynomial $h(r)$ of degree at most $s$, while retaining
the roots of $f^{(D-n)}$. Determine the universal finite-difference
identities for these means, including their polynomial continuation
at $n=0$, and characterize the choices of $h$ for which that
endpoint vanishes for every $f$.

\section*{Task 5: relations involving antiderivatives}
The source also studies means over roots of antiderivatives and
proves independence from their integration constants for the
relevant polynomial evaluations. Building on that result, classify
all universal linear relations for a fixed finite collection of
derivative and antiderivative orders. Determine whether one
interpolating polynomial accounts for the whole collection.

\paragraph{Verification and status.}
The accompanying exact verifier checks generic symbolic
coefficients through degree nine and separately computes means
using traces of polynomial evaluations on companion matrices.
The tests include even degrees, nonmonic polynomials, complex
coefficients, and repeated roots. They supplement the all-degree
proof. Tasks 2--5 are proposed directions and receive no solved
credit. The complete 27-page current primary and a bounded
prior-claim search were checked by the authoring agent; no
independent peer review or proof assistant verification is claimed.

\begin{thebibliography}{9}
\bibitem{W} G. G. Wojnar, D. Sz. Wojnar and L. Q. Brin,
\emph{Universal Peculiar Linear Mean Relationships in All
Polynomials}, arXiv:1706.08381v2 (23 October 2017),
definitions on pages 1--2 and Conjecture 2 on page 9.
The PDF title omits ``Peculiar.''
\url{https://arxiv.org/abs/1706.08381}.
\end{thebibliography}
\end{document}
